% =============================================================================
\section{Results and Discussion}\label{sec:results}
% =============================================================================

All results use the Omicron test split (1~Dec~2021--31~Aug~2022).

\subsection{Forecasting Accuracy}\label{ssec:e1}

\Cref{tab:forecast_accuracy} reports per-horizon RMSE and the
$h{=}28$ underestimation rate against four baselines under a common
$32$-origin rolling-origin protocol. The \emph{raw} PinnGRU point
forecast (frozen PINN encoder of \Cref{ssec:pinn} with the GRU head of
\Cref{ssec:temporal}) wins $h{=}7$ and $h{=}14$---paired origin-level
tests confirm the $h{=}14$ gain over ARIMA ($-3.88$ MV beds, Wilcoxon
$p{=}0.003$)---but degrades at long
horizons, where a regime-shift over-prediction bias (its $h{=}28$
underestimation rate is only $10.3\%$) leaves it behind ARIMA and
seasonal-naive ($+5.53$ MV beds vs ARIMA at $h{=}28$, $p{=}0.006$). A
leakage-free per-origin recalibration (\Cref{ssec:recal}) removes this
bias and blends with persistence at long horizons: the recalibrated
model is best at $h{=}7$, $14$ and $21$---cutting the operational
$14$-day RMSE by $25\%$ over ARIMA and $42\%$ over the per-region
GRU---and at $h{=}28$ ties ARIMA ($28.5$ vs $28.0$), with only
seasonal-naive ($26.0$) still ahead. Its $h{=}28$ underestimation rate
rises to $37\%$: the point forecast is now a well-calibrated median, not
a systematic over-predictor.

A second observation concerns the non-ARIMA models. XGBoost and the
per-region GRU are lag-only learners bounded by the $28$-day window;
unable to extrapolate into Omicron's altered case-to-hospitalisation
regime, they degrade sharply ($h{=}28$ RMSE $58.9$ and $44.0$). ARIMA's
per-origin refit keeps its degradation modest, and the recalibrated
PinnGRU matches it across horizons, consistent with the ``w/o PINN
params'' ablation isolated next.

The bottom block of \Cref{tab:forecast_accuracy} ablates three
load-bearing components. Replacing the pinball loss with MSE inflates
$h{=}14$ RMSE by $29\%$; zeroing the four PINN parameter features
inflates it almost two-fold ($14.46\!\to\!27.25$, the largest regression
in the table); and skipping PINN pre-training inflates it by $19\%$
($\to\!17.27$), confirming that the physics prior, not capacity, drives
the gain over the per-region GRU.

\begin{table}[t]
\caption{Forecast accuracy and ablations, macro-averaged over seven
NHS regions and $32$ rolling origins. \textbf{Bold}: best;
\underline{underline}: second best, per column (among the baselines and
the proposed model). MASE is the horizon-averaged mean absolute scaled
error (scaled by the in-sample seasonal-naive($7$) MAE; $<\!1$ beats
that benchmark, lower is better).}
\label{tab:forecast_accuracy}
\centering
\renewcommand{\arraystretch}{1.02}
\setlength{\tabcolsep}{4pt}
\begin{tabular}{lcccc@{\hskip 1em}c}
\toprule
\textbf{Model} & \multicolumn{4}{c}{\textbf{RMSE by horizon (days)}}
               & \textbf{MASE} \\
\cmidrule(lr){2-5}
               & 7 & 14 & 21 & 28 & \textbf{(all $h$)} \\
\midrule
Seasonal-naive(7)             & 13.53 & 18.16 & \underline{22.31} & \textbf{25.98} & 0.63 \\
ARIMA per region              & 11.33 & 17.80 & 23.41 & \underline{28.02} & \underline{0.61} \\
XGBoost (lag features)        & 12.45 & 22.90 & 36.42 & 58.88 & 0.89 \\
GRU per region                & 14.13 & 22.92 & 33.35 & 43.97 & 0.90 \\
PinnGRU (raw)                 & \underline{11.02} & \underline{14.46} & 24.19 & 34.63 & 0.77 \\
\textbf{PinnGRU+recal (prop.)} & \textbf{10.56} & \textbf{13.34} & \textbf{20.79} & 28.54 & \textbf{0.56} \\
\midrule
\quad w/o cost-asymmetric loss & 14.79 & 18.70 & 27.09 & 41.47 & 0.97 \\
\quad w/o PINN pre-training   & 12.26 & 17.27 & 22.56 & 34.73 & 0.76 \\
\quad w/o PINN params         & 23.97 & 27.25 & 32.94 & 44.59 & 1.23 \\
\bottomrule
\end{tabular}
\end{table}

\subsection{Allocation Performance Across Solvers}
\label{ssec:e2_e3}

\Cref{tab:allocation} compares the allocation policies on the
seven-region problem with $B = 213$~beds ($20\%$ of total baseline), four-hour
mutual-aid cap, and PinnGRU $q^{(0.1)}\!/q^{(0.5)}\!/q^{(0.9)}$
forecasts at the first Omicron-period origin. Every policy's chosen
allocation is re-scored by the common LP re-evaluation over the full
three-scenario set so $E[u]$ and $u^{\text{worst}}$ are reported on
equal footing.

\textbf{(i)} The \emph{robust LP} achieves the lowest transfer
burden in the table ($672$~bed-km, $4.6\times$ below
population-proportional and $1.7\times$ below greedy shortage-first),
but this is a tie-break among policies with identical rounded $E[u]$
and $u^{\text{worst}}$, not a coverage improvement. \textbf{(ii)} The
\emph{deterministic LP} spends only $51$ of the $213$ available beds
and pays $\sim\!3\times$ higher $E[u]$ and $u^{\text{worst}}$,
quantifying the value of $\pi_s$-weighted expectation over median-only
optimisation. \textbf{(iii)} The robust LP therefore wins on
\emph{routing} rather than coverage. Once the budget binds, every
budget-exhausting policy ties on $E[u]$ and $u^{\text{worst}}$ and
differs only on transfer burden.

\begin{table}[t]
\caption{Allocation comparison under $B = 213$~beds, $\tau = 240$~min,
and three PinnGRU scenarios.}
\label{tab:allocation}
\centering
\renewcommand{\arraystretch}{1.02}
\setlength{\tabcolsep}{4pt}
\begin{tabular}{lrrrrr}
\toprule
\textbf{Policy} & \textbf{Beds} & \textbf{$E[u]$}
                & \textbf{$u^{\mathrm{worst}}$} & \textbf{Transfer}
                & \textbf{Time}\\
                & & & & \textbf{(bed\,km)} & \textbf{(s)}\\
\midrule
Population-proportional    & 212.7 & 29.0  & 144.8 & 3119.8  & 0.00 \\
Demand-proportional        & 212.7 & 29.0  & 144.8 & 2615.3  & 0.00 \\
Greedy shortage-first      & 212.7 & 29.0  & 144.8 & 1140.0  & 9.98 \\
Deterministic LP           &  51.4 & 83.3  & 416.4 & 1175.1  & 0.03 \\
\textbf{Robust LP ($\lambda_3{=}1$)}
                           & 212.7 & 29.0  & 144.8 & \textbf{671.7} & 0.04 \\
\bottomrule
\end{tabular}
\end{table}

Inspecting the regional structure behind \Cref{tab:allocation}
(\Cref{fig:alloc_budget}a), the deterministic LP surges only Midlands
and South East (where median demand exceeds baseline); budget-exhausting
policies all cover those two and differ on the rest, the robust LP
finding a compromise with non-trivial North~West and London surge.

\begin{figure}[ht]
\centering
\includegraphics[width=\textwidth]{fig_alloc_budget.png}
\caption{(a)~Per-region peak surge allocation (beds) by policy at the
first Omicron origin. (b)~Exact robust-LP cost--shortage frontier
showing how expected and worst-case unmet demand fall as the surge
budget $B/\sum_r C_r$ grows. This is the decision-relevant trade-off a
planner navigates.}
\label{fig:alloc_budget}
\end{figure}

\subsection{Forecast-Quality and Parameter Sensitivity}
\label{ssec:e5_e6}

\paragraph{Forecast-quality robustness.}
Substituting each forecaster into the robust LP at the $29$~Dec~$2021$
origin, the $q^{(0.9)}$ peak sums (PinnGRU $1421$, XGBoost $1105$, GRU
$973$, ARIMA $875$, seasonal-naive $745$~beds) all exceed the realised
peak ($797$) except seasonal-naive. Realised unmet is zero throughout,
including the oracle, so this origin cannot separate policies on
realised coverage: forecast quality buys \emph{efficiency} at the
$q^{(0.9)}$ margin, not coverage, in this mild Omicron instance.

Across all $32$ Omicron origins the pattern holds: robust LP,
population-proportional, and demand-proportional allocations share a
mean expected unmet of $2.9$, but the robust LP cuts mean transfer
burden to $76$~bed-km versus $406$ and $337$~bed-km. The proportional
baselines fill the full $212.7$-bed budget at every origin, whereas the
robust LP scales surge to local demand (mean total surge $35.4$~beds,
hitting the cap only at peaks), so the routing-efficiency gain holds
across the full demand distribution.

\paragraph{Parameter sensitivity.}
A one-at-a-time sweep isolates each control. The budget $B$ dominates
(\Cref{fig:alloc_budget}b): $B/\sum_r C_r$ from $0.20$ to $0.10$ doubles
$E[u]$ and triples $u^{\text{worst}}$, while $0.30$ cuts $E[u]$ to
$\sim\!1/4$. $\lambda_3$ is \emph{non-binding} at $B=0.20$ (the
expectation-minimising solution already clears the high tail); only a
tighter budget ($\leq\!15\%$) re-engages it. The travel cap
$\tau\!=\!120$~min cuts transfer burden by $34\%$ at a $3\%$ rise in
$E[u]$; $\tau\!\geq\!240$~min is slack.

% Table 3 (one-at-a-time sensitivity sweep) removed for the 12-page UKCI
% limit: the budget/tail-penalty/travel-cap effects are reported in the prose
% above and the exact budget frontier is Fig.~4(b). The full sweep CSVs
% (e6_*.csv) remain in the repository for the journal extension.
